\documentclass[11pt]{article}
\usepackage[utf8]{inputenc}
\usepackage[T1]{fontenc}
\usepackage[english]{babel}
\usepackage{graphicx,amsmath,hyperref,booktabs}
% preamble comment with a "quote" and {braces
\newcommand{\pkg}[1]{\texttt{#1}}
\title{Lumen Bakery \emph{Handbook}}
\author{The Lumen Team}
\date{\today}
\begin{document}
\maketitle
\tableofcontents

\begin{abstract}
This handbook explains how the bakery works, from the first oven to the last customer.
\end{abstract}

\section{Welcome}\label{sec:welcome}
Fresh bread is ready at \textbf{six} every morning---see Section~\ref{sec:hours} and \cite[p.~3]{baker99}.
Prices start at 3\,\% of the \$12 loaf; ask about R\&D loaves, a\_b and {\bf grouped} words.
Inline math $E = mc^2$ and \(a+b\) stay, and so does display math:
\[ \sum_{i=1}^n i = \frac{n(n+1)}{2} \]
$$ x^2 + y^2 = z^2 $$

\subsection*{Hours}\label{sec:hours}
\begin{itemize}
  \item Open from six to noon.
  \item[Sundays] Closed\footnote{Except on holidays; call us first.}.
  \item Visit \url{https://lumen.example/hours%20today} or \href{https://lumen.example}{our website}.
\end{itemize}

\begin{enumerate}
  \item Mix the dough. % remove for the demo
  \item Let it rest.
\end{enumerate}

\begin{equation}
  a^2 + b^2 = c^2 \text{ never translated}
\end{equation}
\begin{align*}
  f(x) &= x^2 \\
  g(x) &= x^3
\end{align*}

\begin{verbatim}
printf("Do not translate this: 100%");
\end{verbatim}
Use \verb|code_here| and \texttt{plain text} in a sentence.

\begin{figure}[htbp]
  \centering
  \includegraphics[width=0.5\textwidth]{bread.png}
  \caption[Bread]{A fresh loaf of bread.}
  \label{fig:bread}
\end{figure}

\begin{table}[h]
\centering
\begin{tabular}{|l|r|}
  \hline
  Product & Price \\
  \hline
  Rye bread & 3.50 \\
  Sourdough & 4.20 \\ \hline
\end{tabular}
\caption{Today's prices}
\end{table}

\begin{lstlisting}[language=Python]
print("no")
\end{lstlisting}

Written with \LaTeX\ and plain \TeX. See \S 2 and \ldots{} more.
\input{appendix}
\bibliographystyle{plain}
\bibliography{refs}
\end{document}
Text after the document end is still read.
